% part: intuitionistic-logic
% chapter: propositions-as-types

\documentclass[../../../include/open-logic-chapter]{subfiles}

\begin{document}

\olchapter{int}{pty}{Proposisi sebagai Tipe}

\begin{editorial}
  Ini merupakan draf bab tentang korespondensi Curry--Howard yang
  \emph{sangat eksperimental}. Bab ini memerlukan lebih banyak penjelasan
  dan motivasi, dan kemungkinan memuat kesalahan serta penghilangan.
  Bukti normalisasi harus ditinjau dan diperluas. Tidak ada contoh untuk
  tipe produk. Konversi permutasi dan penyederhanaan belum dibahas. Bab
  ini akan jauh lebih masuk akal setelah tersedia pula materi tentang
  kalkulus lambda (bertipe), yang pada dasarnya dipraanggapkan di sini.
  Gunakan dengan sangat hati-hati.
\end{editorial}

\olimport{introduction}
%\olimport{sequent-natural-deduction}
\olimport{proof-terms}
\olimport{proofs-to-terms}
\olimport{terms-to-proofs}
\olimport{types}
\olimport{reduction}
\olimport{type-preservation}
\olimport{normalization}

\OLEndChapterHook

\end{document}
